\section{Introduction}
\label{sec:intro}

A managed agent platform runs an agent as a cloud workload with an identity of
its own. On \agentcore an agent runtime is a resource with an execution role,
and the account that holds one agent usually holds several. They share a
region, a network, a logging configuration, and often a role.

Tools that audit cloud permissions ask a question about one identity: given
this principal, what can it reach. The question is answered against the
principal's own policies, the resource policies it meets, and the roles it may
assume. It is the right question for a human user or a service account, and
every finding it produces is about that one identity.

It does not express the problem in this paper. When two agents live in
one account, the interesting failure is not that agent A can reach too much on
its own. It is that agent A can reach agent B, change what agent B runs, and
have that code execute under agent B's role. Nothing in A's own permission set
says it can reach B's data. The escalation exists in the composition of two
identities, and a tool that only ever looks at one of them will report both as
clean.

The paper writes that as a predicate over two sets: what the deployment
authorised a principal to reach, and what it reaches once every edge is
followed, including edges that run through another agent. Cross-agent
escalation is the difference between them.
Section~\ref{sec:model} states both sets and the predicate, and
Section~\ref{sec:approach} gives the rules that make the second computable
from infrastructure as code.

\subsection{Two routes to a sibling agent's role}

There are two ways to act with a sibling agent's role. The first reads its
credentials: a single permission runs a command inside the agent's own
sandbox, and the platform's metadata service hands the role over. That route
is published~\cite{garcia2026agentcorepaths}, needs no
\iam{iam:PassRole}, and leaves the deployed agent exactly as it was.

The second replaces what the agent runs.
\api{Update\allowbreak AgentRuntime} marks two
fields as required: the artifact holding the agent's code, and the ARN of the
role it runs as. One call therefore decides what runs, who it runs as, and,
through the runtime's environment variables, what it is aimed at. Any caller
that replaces an agent's artifact must also pass a role, so this route always
crosses \iam{iam:PassRole}. 
This paper takes the second because its consequence outlives the call: the
agent keeps its ARN, its endpoint and its callers, and answers them with
attacker code, which Section~\ref{sec:case-study} measures by stripping the
attacker of every permission first. The command route ends when the credential
is rotated. A replaced artifact does not.

Cross-agent escalation on this platform has been reported.
Section~\ref{sec:industry} sets out what those accounts cover and what
they do not: the path that replaces an existing runtime rather than creating a
new one, the experiment that isolates the role name as the control, and any
measurement of how often the preconditions occur.

\subsection{Contributions}

\begin{itemize}
\item A predicate for cross-agent privilege escalation
  (Equation~\ref{eq:esc}) and a delegation edge rule that makes it computable
  from infrastructure as code, implemented in \tool
  (Sections~\ref{sec:approach} and~\ref{sec:implementation}).
\item A measurement over \nProjects public Terraform projects that deploy
  agent runtimes and the CloudFormation and CDK deployments in the same
  collection (Sections~\ref{sec:methodology} and~\ref{sec:results}), finding
  \nChainsFound complete cross-agent escalations, with the bound that supports
  and the preconditions that are common. A \foundry corpus of \nFoundryDirs
  template directories is surveyed separately
  (Section~\ref{sec:foundry}).
\item A live demonstration in an AWS account showing that an attacker holding
  one AWS managed policy can take over a sibling agent, and that the gate is a
  name match in that policy's \iam{iam:PassRole} condition
  (Section~\ref{sec:case-study}).
\item A claim by claim account of what was verified and what was not
  (Appendix~\ref{sec:claims}), and the artifacts to rebuild every number
  (Section~\ref{sec:availability}).
\end{itemize}

